\documentclass[10pt,a4paper]{amsart}
\usepackage[margin=19mm]{geometry}
\usepackage{amsmath,amssymb,mathtools}
\usepackage[scr=rsfs,cal=euler]{mathalfa}
\usepackage{xfrac}
\usepackage[unicode,hidelinks,bookmarks=false]{hyperref}
\newcommand{\ie}{i.e.\ }
\newcommand{\cf}{cf.\ }
\newcommand{\resp}{respectively\ }
\newcommand{\Subsection}{\subsection}
\providecommand{\nobreakdash}{\nobreak\hbox{-}}
\setlength{\emergencystretch}{2em}
\multlinegap0pt
\usepackage{amssymb}
\usepackage[shortlabels]{enumitem}
\usepackage{tikz}
\usepackage{algorithm}\usepackage{algpseudocode}
\newtheorem{theorem}{Theorem}
\newcommand\Me{\mathcal{M}^\mathrm{ext}}
\newcommand\ba{\backslash}
\title{Extensions and Deletions of matroid classes closed under flats}
\author{Jagdeep Singh, Vaidy Sivaraman}
\date{Source: arXiv:2403.15496v3 (CC BY 4.0)}
\begin{document}
\maketitle

\noindent\emph{Source and licence.} Extensions and Deletions of matroid classes closed under flats. Version arXiv:2403.15496v3 (CC BY 4.0), distributed by arXiv under the Creative Commons Attribution 4.0 International licence, http://creativecommons.org/licenses/by/4.0/, as recorded in the arXiv licence metadata for this exact version and shown on the abstract page https://arxiv.org/abs/2403.15496v3. Full licence notice: \texttt{licenses/arxiv-2403.15496-CC-BY-4.0.txt}.\par\smallskip

\noindent\textit{Editorial scope.} Counted unit: the article's theorem bound\_k\_general (source0.tex lines 168--171). The article's complete original proof of it is delivered with it (source0.tex lines 174--180, one \texttt{proof} environment of 7 lines). No mathematics is written, repaired or re-derived: the statement and the proof are the article's own lines, in the article's own notation, and no retained line is substituted or re-flowed.  No further source-local context is needed: every cross-reference of every retained line resolves inside the two delivered spans.  Nothing was omitted from any retained span, so the package carries no omission record.  No retained line contains a citation, so the wrapper carries no bibliography and no bibliography entry.  No retained line contains a figure environment, an image-inclusion command or drawing code, so no image is delivered and the package declares no asset dependency.  Editorial formatting only, in addition: an AMS \texttt{amsart} preamble that restates the article's own declarations and macros used by the retained lines, namely the environments \texttt{theorem} on separate counters, together with the standard packages those lines need.  The retained environments print in this document as Theorem 1 in order of appearance, whereas the article prints the same items with the numbering of its own full document.  The complete native source of the article as distributed in the arXiv e-print for this exact version is bundled unchanged as \texttt{sources/156500/source0.tex}.
\begin{theorem}
\label{bound_k_general}
Let $\mathcal{M}$ be a hereditary class of matroids, and let $\mathcal{F}$ be the set of its forbidden flats. Let $r$ denote the maximum rank of a flat in $\mathcal{F}$. If $M$ is a forbidden flat for $\Me$, then the rank of $M$ is at most $\max\limits_{F \in \mathcal{F}}{\{2r, r(F) + |F|(r-1)\}}$. 
\end{theorem}

\begin{proof}
Note that $M$ is not in $\mathcal{M}$ so $E(M)$ has a flat $F_0$ such that the restriction $M|F_0$ is a forbidden flat for $\mathcal{M}$. Suppose that $F_0 = \{e_1, \ldots, e_k\}$ and $r(F_0) = s$. Since for each $i$ in $\{1, \ldots, k\}$, the matroid $M \ba e_i$ is not in $\mathcal{M}$, we have a subset $F_i$ of $E(M)-e_i$ such that the restriction $(M \ba e_i)|F_i$ is a forbidden flat for $\mathcal{M}$. Observe that, if $\bigcup_{i=0}^{k} F_i$ is contained in a hyperplane $H$ of $M$, then the matroid $M|H$ is not in $\Me$, a contradiction to the minimality of $M$. Therefore, $r(M) = r(\bigcup_{i=0}^{k} F_i)$.

First suppose that the intersection of a pair in $\{F_0, F_1, \ldots, F_k\}$, say $F_m$ and $F_n$, is empty. Then $r(M)=r(F_m \cup F_n)$. If not, then $F_m \cup F_n$ is contained in a hyperplane $H$ of $M$, and $M|H$ is not in $\Me$, a contradiction. Since $r(F_m) \leq r$ and $r(F_n) \leq r$, it follows that $r(M) = r(F_m\cup F_n) \leq 2r$. 

Therefore we may assume that each such pair has a non-empty intersection. It now follows from the submodularity of the rank function that $r(\bigcup_{i=0}^{k} F_i) \leq s + k(r-1)$. Since $s = r(F_0)$ and $k = |F_0|$ for the forbidden flat $F_0 \in \mathcal{F}$, we conclude that $r(M) \leq \max\limits_{F \in \mathcal{F}} \{2r, r(F) + |F|(r-1)\}$. This completes the proof.
\end{proof}

\end{document}
